\section{Proofs for negative-premise transition specifications}

\begin{annexitem}{annex:p09-np-positive}{thm:p09-np-positive}{Regular-cardinal closure and genuine well-founded proofs}
\begin{proof}
For positive ground rules define \(F(X)\) as the heads whose full
premise set lies in \(X\). It is monotone. Starting at empty,
take \(X_{\xi+1}=F(X_\xi)\) and unions at limits. The chain
is increasing: use monotonicity at successors; at a limit, each
earlier generated head still belongs to \(F\) of the union.
If all premises of a rule lie in \(X_\lambda\), assign to each
a stage below \(\lambda\). There are fewer than \(\lambda\)
such stages, so regularity bounds them by one \(\eta<\lambda\).
The head lies in \(X_{\eta+1}\subseteq X_\lambda\). Thus
\(F(X_\lambda)=X_\lambda\). Every model contains every stage
by transfinite induction, hence this is the least model; fixed-point
equality also supplies a supporting root rule for every member.
A set of rules admits a common infinite regular bound strictly
above the supremum of their premise cardinalities in the stated
classical set theory. Ground instantiation cannot increase premise
cardinality. Finite premises permit \(\omega\).

Induct on first generation stage to build a proof for every generated
head: choose a generating ground rule and proofs for its premises
from earlier stages. Along a branch the chosen stages strictly
decrease, excluding an infinite branch; the children form a set.
Conversely well-founded induction on any positive proof puts all
its children, then its root, in the least model. This is not merely
finite-height induction: infinitely branching trees are included.
\end{proof}
\end{annexitem}

\begin{annexitem}{annex:p09-np-stable}{thm:p09-np-stable}{Reduct proofs and negative-assumption proofs coincide}
\begin{proof}
A positive proof in \(T^M\) uses ground rules whose erased negatives
hold in \(M\). Restore those negatives as leaves and collect their
labels in a set \(N\). Adding leaves does not create an infinite
branch, so \(T\vdash N/\alpha\) and \(M\models N\).
Conversely prune the negative hypothesis leaves from such a proof.
Every retained node is a rule of \(T^M\), giving a positive proof.
The previous lemma identifies its conclusions with \(\Gamma_T(M)\).
If \(M=\Gamma_T(M)\), any ground rule enabled in \(M\) is
retained in its reduct and has its positives in that least positive
model. Closure gives its head, proving modelhood. Restored proof
leaves establish well-supportedness; a proof's root gives ordinary
support. Conversely, well-supportedness gives
\(M\subseteq\Gamma_T(M)\), while modelhood makes \(M\) a
model of the reduct and gives \(\Gamma_T(M)\subseteq M\).
Thus the two meanings coincide. The positive loop below refutes
the converse from merely supported modelhood.
\end{proof}
\end{annexitem}

\begin{annexitem}{annex:p09-np-well-founded}{thm:p09-np-well-founded}{The least alternating pair and its exact information order}
\begin{proof}
For \(M\subseteq N\), every negative true in \(N\) is true
in \(M\). Hence \(T^N\subseteq T^M\); positive proof
inclusion gives \(\Gamma_T(N)\subseteq\Gamma_T(M)\).
Its square \(G\) is therefore monotone.
For any monotone powerset operator, let \(L\) be the intersection
of all \(X\) satisfying \(G(X)\subseteq X\). Monotonicity
gives \(G(L)\subseteq L\), and then
\(G(G(L))\subseteq G(L)\). Thus \(G(L)\) is itself in
the intersected family, so \(L\subseteq G(L)\). Equality
proves least fixed-point existence. Transfinite iteration from
empty, with unions at limits, gives the same value: it is increasing
and bounded by every pre-fixed set, and cannot strictly add a new
atom at every stage below \(|B|^+\). Choosing the least fresh
atom in a fixed well-order of \(B\) would otherwise inject that
larger cardinal into \(B\). No continuity at \(\omega\) is assumed.

For \(L_*=\mu G\), \(\Gamma_T(L_*)\) is another fixed
point of \(G\), since the square commutes with \(\Gamma_T\).
Leastness gives \(L_*\subseteq U_*:=\Gamma_T(L_*)\), and
\(\Gamma_T(U_*)=L_*\). For any other three-valued stable
pair, \(G(L)=L\) implies \(L_*\subseteq L\); antitonicity
gives \(U=\Gamma_T(L)\subseteq U_*\). A stable model is
the pair \((M,M)\), hence lies between the bounds. If the
bounds agree their common value is the unique stable model.
The fully calculated example \(E\) below refutes the converse.
\end{proof}
\end{annexitem}

\begin{annexitem}{annex:p09-np-proof-meanings}{thm:p09-np-proof-meanings}{Universal negative blocking and completeness of the bounds}
\begin{proof}
Induct on a supported proof against a supported model \(M\).
Positive nodes use enabled ground rules, so are in \(M\).
If a negative node had a denier in \(M\), support would supply
a rule for that denier with premises all true in \(M\). Its
blocking condition finds one of those premises denying a child
already true by induction, a contradiction. This proves soundness.

For a \(ws\)-proof and a three-valued stable pair \(L\subseteq U\),
induct simultaneously that positive nodes belong to \(L\), and
negatives have no denier in \(U\). A positive rule is retained
in \(T^U\) because its negative children have no deniers there;
its positives are in \(L=\Gamma_T(U)\), giving its head by
closure. If a negative node \(\eta\) had a denier
\(\gamma\in U=\Gamma_T(L)\), the previous lemma gives
\(T\vdash N/\gamma\) with negative \(N\) true in \(L\).
The blocking condition supplies a member of \(N\) denying a
child. That child must be positive, hence in \(L\) by induction,
contradicting the negative hypothesis. Applying this soundness to
the existing least pair proves consistency of \(ws\)-provability.

For the converse, iterate \(G=\Gamma_T^2\) from empty as
\(L_\xi\), and set \(U_\xi=\Gamma_T(L_\xi)\). Induct
that every positive in \(L_\xi\) has a \(ws\)-proof.
Once this holds, every negative \(\eta\) with no denier in
\(U_\xi\) has a \(ws\)-proof using all positives of
\(L_\xi\) as children. For any ordinary \(N/\gamma\)
denying \(\eta\), the fact
\(\gamma\notin\Gamma_T(L_\xi)\) means that \(N\) is
not true in \(L_\xi\). A member of \(N\) has a denier
in \(L_\xi\), exactly the required child. The child set and
chosen well-founded proofs are sets; grafting them under one root
is well-founded. At stage 0 this is the empty-child, vacuous-blocking
case. At a successor, every
\(\alpha\in\Gamma_T(U_\xi)=L_{\xi+1}\) has an ordinary
proof under negative hypotheses true in \(U_\xi\); replace
those leaves by the just constructed negative proofs. Any branch
follows one finite original branch and then one well-founded graft,
so the resulting proof is well-founded. At a limit positives came
from earlier stages. At stationarity this proves all positives in
\(L_*\) and all negatives with no denier in \(U_*\).
Soundness gives the reverse inclusions. Singleton-denier negatives
exist for each positive atom in the extended proof language;
therefore \(ws\)-completeness is exactly \(L_*=U_*\).
\end{proof}
\end{annexitem}

\begin{annexitem}{annex:p09-np-stratified}{thm:p09-np-stratified}{The full ordinal supported/stable construction}
\begin{proof}
Outer recursion is along the set ordinal \(\kappa\). At rank
\(\beta\), all negative deniers are strictly lower, so their
truth is decided by \(I_{<\beta}\). Positive premises are lower
or same rank. Discard failed lower tests and erase true lower
premises. The remaining positive rules on \(B_\beta\) have
fewer than \(\lambda\) premises. The positive lemma proves
that the displayed union is the least fixed point of \(F_\beta\).
This proves well-definedness and the regular-cardinal bound.
For a ground rule enabled in \(M_S\), its positive premises
are lower or in \(I_\beta\), and its negatives hold exactly
in the lower part because all deniers have lower rank. Closure
gives its head. Conversely, fixed-point equality provides a ground
rule witness for each member. Thus \(M_S\) is a supported model.

Being a model makes it a model of \(T^{M_S}\), hence
\(\Gamma_T(M_S)\subseteq M_S\). Induct on stratum and,
within each, generation stage to prove the reverse inclusion.
Each generating rule is retained in \(T^{M_S}\); its lower
positives have entered that least model by outer induction, and
its same-rank positives by inner induction. Its head therefore
enters by closure. Limits use unions. Thus \(M_S\) is stable.

Here is the needed localization fact, with proof. For any positive
rule set whose premise ranks never exceed head rank, its global
least model \(Q\) at rank \(\beta\) equals the least local
closure with the actual \(Q_{<\beta}\) fixed. Local-to-global
inclusion follows by stage induction because \(Q\) is a model.
For the other inclusion, take a positive proof of a rank-\(\beta\)
head. All nodes have rank at most \(\beta\). Treat lower-rank
nodes as known members of \(Q_{<\beta}\); well-founded induction
on the remaining same-rank proof puts its nodes in the local least
closure. This proves equality, not an assumed commutation of closures.

For another stable \(N\), induct on rank that its lower part
equals that of \(M_S\). Negative tests at the current head rank
then agree, and localization of \(N=\Gamma_T(N)\) yields
exactly the same local least closure \(I_\beta\). Thus stable
uniqueness holds. The stable definition mentions neither a rank
map nor a cardinal bound, so uniqueness gives independence of both.
Apply localization likewise to
\(L_*=\Gamma_T(U_*), U_*=\Gamma_T(L_*)\). At rank
\(\beta\), induction makes their lower parts both equal to
\((M_S)_{<\beta}\). Their negative tests and positive local
frontiers therefore coincide, giving \(I_\beta\) in both bounds.
Consequently \(L_*=U_*=M_S\), and the proof-meaning theorem
gives \(ws\)-completeness.

Under strict stratification every positive premise is also lower.
The local operator is then constant. In a supported model, a head
is present exactly when some ground rule's lower premises hold.
Rank induction forces that constant value at every stratum, giving
unique supported modelhood. Ordinary stratification alone cannot
give this last conclusion, as the source positive self-loop shows.
\end{proof}
\end{annexitem}

\begin{annexitem}{annex:p09-np-countermodels}{thm:p09-np-countermodels}{The six source examples and two precise further countercases}
\begin{proof}
On the stated predicate bases, \(T_1\)'s models are both singletons
and the pair; its supported models are the singletons. Its reduct
operator is \(\{h:k\notin M\}\cup\{k:h\notin M\}\),
with exactly those two fixed points. They are incomparable minimal
models, not two different least models; there is no least model.
For \(T_2\), empty fails its enabled rule. The singleton is a
model but cannot support its own head because its only rule is
disabled. Its reduct alternates empty and the singleton, with no
fixed point. For \(T_3\), both singletons and the pair are models,
but only \(\{h\}\) is supported, since no rule has head \(k\).
Its reduct is \(\{h\}\) when \(k\notin M\), and empty
otherwise, so its only stable model is \(\{h\}\).

For \(T_4\), the models are \(\{h\}\) and \(\{h,k\}\).
The former cannot support \(h\); the latter supports \(h\)
from \(k\) and \(k\) from itself. In every reduct those two
positive rules have no grounded axiom and generate nothing. The
remaining negative rule gives a fact \(h\) exactly when
\(h\notin M\), so there is no stable fixed point. For \(T_5\)
empty and the singleton both satisfy support/model conditions,
but no well-founded positive proof can derive \(h\): it would
have an infinite repeated branch. Every reduct least model is
empty. For \(T_6\), the two supported models are \(\{h\}\)
and \(\{k\}\); the pair cannot support \(k\). Its reduct
generates \(k\) exactly when \(h\notin M\), and never
generates \(h\). Its only stable model is \(\{k\}\).
Squaring these operators gives the claimed total well-founded
values for \(T_3,T_5,T_6\). For \(T_1,T_2,T_4\), the
lower bound stays empty, with upper bounds respectively
\(\{h,k\},\{h\},\{h\}\).

In \(E\), a stable model containing \(h\) must omit \(k\);
the \(r\)-rule is disabled, giving exactly \(\{h\}\).
Omitting \(h\) forces \(k\), but requires \(r\) exactly
when it is absent, impossible. Thus stability is unique. However,
\(\Gamma_E(\varnothing)=\{h,k,r\}\) and
\(\Gamma_E(\{h,k,r\})=\varnothing\). The least squared
fixed point is empty and its upper bound is all three, so unique
stability does not imply total well-founded semantics.

For \(J\), stage 1 contains label 0, and stage \(n+1\) contains
label \(n\). Their union at \(\omega\) contains all numerical
labels but not \(b\). The infinitary rule is then enabled, so
this union is not a model; \(b\) appears at \(\omega+1\).
Its proof has arbitrarily deep finite branches, no infinite branch.
A regular bound strictly larger than its countable premise set,
for example \(\omega_1\), includes the necessary stage. This
is a genuine transition-rule obstruction to assuming omega closure.
\end{proof}
\end{annexitem}

\begin{annexitem}{annex:p09-np-processes}{thm:p09-np-processes}{Every Table~1--3 rank case and actual process transition}
\begin{proof}
All nine BPA rules are positive, so rank 0 is a stratification.
For priority, count occurrences of \(\theta\) in the source
of every action or predicate judgement. At a choice or product,
the count is the sum of the child counts, so every premise source
has no greater rank. This checks the four choice and three product
rules, including the two-premise success and right-action cases.
The two axioms have no tests. At a priority root the rank is one
plus the child's rank. Its success/action positives and every
possible denier of each higher-action negative have that strictly
lower child rank, independently of target. This also covers an
infinite set of higher actions.

For the separate timed TSS, ordinary actions and success have rank 0;
timed judgements have source choice-count. BPA rules still meet
their rank conditions. The delay axiom has no premises. At a
product, choice-count is the sum of child counts, bounding left
and right timed premises and the success test. At a choice it is
one plus that sum, strictly above both child counts. This checks
both positives in the both-delay rule, and all possible negative
deniers in each one-delay rule. These exhaust the six timed rules.
Node-size, independent of target and label, is separately a strict
stratification: every nonaxiom positive premise and every denier
has source a proper immediate subterm of the conclusion source.
This proves unique supported models as well, without calling the
source's rank 0/counting maps themselves strict. No combined
priority/time/recursive system is asserted.

The base process \(a+b\) has exactly its two action transitions
to \(\epsilon\), and no success predicate. With \(a<b\),
the priority \(a\)-rule fails its higher \(b\) guard; the
\(b\)-rule has no larger label and gives \(\theta(\epsilon)\).
No other rule gives a priority-headed action conclusion, so the
calculation is exhaustive. Its target succeeds by the success rule.
Each delayed atom has its sole time transition to that atom;
an undelayed atom has no timed rule. Thus \(\sigma_d(a)+b\)
uses exactly the left-only rule and goes to \(a\). When both
summands are delayed, the both-delay rule gives \(a+b\); each
one-sided negative guard fails. These are exact source-rule
endpoints, not declarations of absence after bounded search.
\end{proof}
\end{annexitem}
